\documentclass[10pt, final, journal, letterpaper, twocolumn]{IEEEtran}

% Fonts and mathematics
\usepackage{amsmath,amssymb}
\usepackage[T1]{fontenc}
\usepackage{newtxtext,newtxmath}

% Hyperlinks and author identifiers (hyperref before glossaries and cleveref)
\usepackage[hidelinks]{hyperref}
\usepackage{orcidlink}

% Figures, tables, and captions
\usepackage{graphicx}
\usepackage{booktabs}
\usepackage{xcolor}
\usepackage[font=small]{caption}

% Acronyms (abbr.tex is the shared list; entries specific to this paper follow)
\usepackage[acronym]{glossaries}
\loadglsentries{abbr.tex}
\newacronym{rrs}{RRS}{radio resource scheduling}
\newacronym{nmse}{NMSE}{normalized mean square error}
\newacronym{eaxc}{eAxC}{extended antenna-carrier}
\newacronym{oru}{O-RU}{O-RAN radio unit}
\newacronym{odu}{O-DU}{O-RAN distributed unit}
\newacronym{prg}{PRG}{precoding resource block group}
\newacronym{dmrs}{DMRS}{demodulation reference signal}
\newacronym{resel}{RE}{resource element}
\glsdisablehyper

% Theorem-like environments
\newtheorem{proposition}{Proposition}
\newtheorem{remark}{Remark}

% Cross-references
\usepackage[capitalise]{cleveref}

% Bibliography
\usepackage[backend=biber, style=ieee, maxnames=2, minnames=1, doi=false, url=false, eprint=false, isbn=false]{biblatex}
\addbibresource{IEEEabrv.bib}
\addbibresource{main.bib}
\AtBeginBibliography{\footnotesize}

% Draft label (remove before submission)
\newcommand{\draftversion}{\textcolor{red}{V0.1}}
% To-do notes (remove before submission)
\newcommand{\todo}[1]{\textcolor{red}{[\textbf{TODO:} #1]}}
\newcommand{\note}[1]{\textcolor{blue}{[Note: #1]}}

% Notation
\newcommand{\Nset}{\mathcal{N}}
\newcommand{\Kset}{\mathcal{K}}
\newcommand{\Uset}{\mathcal{U}}
\newcommand{\Bset}{\mathcal{B}}
\newcommand{\Sru}{\mathcal{S}^{\mathrm{RU}}}
\newcommand{\Sant}{\mathcal{S}^{\mathrm{ant}}}
\newcommand{\Pmax}{P_{\max}}
\newcommand{\SINR}{\mathrm{SINR}}
\newcommand{\SE}{\mathrm{SE}}
\newcommand{\CN}{\mathcal{CN}}
\newcommand{\E}{\mathbb{E}}
\newcommand{\bh}{\mathbf{h}}
\newcommand{\bhh}{\hat{\mathbf{h}}}
\newcommand{\bw}{\mathbf{w}}
\newcommand{\bwh}{\hat{\mathbf{w}}}
\newcommand{\bx}{\mathbf{x}}
\newcommand{\bxh}{\hat{\mathbf{x}}}
\newcommand{\bH}{\mathbf{H}}
\newcommand{\bHh}{\hat{\mathbf{H}}}
\newcommand{\bW}{\mathbf{W}}
\newcommand{\bWh}{\hat{\mathbf{W}}}
\newcommand{\bWt}{\tilde{\mathbf{W}}}
\newcommand{\bE}{\mathbf{E}}
\newcommand{\bEh}{\hat{\mathbf{E}}}
\newcommand{\sone}{7\text{-}1}
\newcommand{\stwo}{7\text{-}2}
\newcommand{\ind}{\mathbb{1}}

\begin{document}

% Working title (to be decided)
\title{\draftversion \quad Minimizing the Fronthaul Load of Cell-Free mMIMO with: Antenna-Level Scheduling, Functional Split, and Compression}

\author{
    \IEEEauthorblockN{Tianzheng Miao\orcidlink{0009-0006-3344-7222}, Md Arifur Rahman\orcidlink{0000-0003-3861-9290}, François Rottenberg\orcidlink{0000-0002-2150-8511}}

     \thanks{
      This work has received funding from the European Union's Horizon 2022 research and innovation program under Grant Agreement No 101120332 (EMPOWER-6G project).
      T. Miao and M. A. Rahman are with the Research and Innovation Department, IS-Wireless, Piaseczno, Poland (e-mail: \{t.miao, a.rahman\}@is-wireless.com). T. Miao and F. Rottenberg are with ESAT-DRAMCO, KU Leuven, Ghent, Belgium (e-mail: \{tianzheng.miao, francois.rottenberg\}@kuleuven.be). 
    }
}

\maketitle

\section{System Model and Problem Formulation}\label{sec:system}

We consider the \gls{dl} of a \gls{cf} \gls{mmimo} network in an \gls{oran} deployment, as in~\cite{girycki2024eucnc}. $N$ geographically distributed \glspl{oru} jointly serve $U$ single-antenna \glspl{ue}, and an \gls{odu}, which acts as the central processing unit of the \gls{cf} network, carries out the \gls{rrs} and the precoding~\cite{ngo2017cellfree,interdonato2019ubiquitous,polese2023oran}. Every \gls{oru} has $K$ co-located antennas, so the network has $M = NK$ antennas. We keep the notation of~\cite{girycki2024eucnc}, where $M$ is the total number of antennas and $K$ the number of antennas per RU; in~\cite{girycki2024eucnc}, however, the $K$ antennas of an RU are spread over the area, whereas here they are co-located and share the large-scale fading.

The problem is that of~\cite{girycki2024eucnc}: select the beamforming type and the functional split that minimize the \gls{fh} load while the sum \gls{se} stays above a fraction $\kappa$ of its maximum. We extend it in three directions. First, the serving antennas are selected on each \gls{prb} with the precoding-aware \gls{rrs} of~\cite{girycki2024twc}, generalized from the selection of single-antenna \glspl{ap} to the activation of individual antennas inside each \gls{oru}. Second, the \gls{fh} payloads are compressed with the methods of the \gls{oran} specification~\cite{oran_cus_v22}, which changes both the \gls{fh} load and the \gls{sinr}. Third, the \gls{odu} schedules with estimated channels. For $K = 1$, uncompressed \gls{fh}, and perfect \gls{csi}, the model reduces to that of~\cite{girycki2024eucnc}.

\textit{Notation:} Boldface lowercase and uppercase letters denote vectors and matrices, respectively. $(\cdot)^*$, $(\cdot)^T$, and $(\cdot)^H$ denote the conjugate, transpose, and conjugate transpose. $[\mathbf{x}]_k$ is the $k$th entry of $\mathbf{x}$ and $[\mathbf{X}]_{i,j}$ the $(i,j)$th entry of $\mathbf{X}$. $\|\cdot\|$ is the Euclidean norm, $|\mathcal{S}|$ the cardinality of a set, $\ind\{\cdot\}$ the indicator function, $\mathbf{I}_K$ the $K \times K$ identity matrix, $\E\{\cdot\}$ the expectation, and $\CN(\mathbf{0}, \mathbf{R})$ the circularly symmetric complex Gaussian distribution with covariance $\mathbf{R}$.

\subsection{Network and Channel Model}\label{sec:channel}

Let $\Nset = \{1, \dots, N\}$, $\Kset = \{1, \dots, K\}$, and $\Uset = \{1, \dots, U\}$ denote the sets of \glspl{oru}, antennas per \gls{oru}, and \glspl{ue}, and let $\Bset$ denote the set of $B$ \glspl{prb} of the carrier. Antenna $k$ of \gls{oru} $n$ is written as $(n, k)$. As in~\cite{girycki2024twc}, a \gls{prb} consists of 12 consecutive subcarriers that experience the same channel. The channel between \gls{oru} $n$ and \gls{ue} $u$ on \gls{prb} $b$ is
\begin{equation}
  \bh_{unb} = \sqrt{\beta_{un}}\, \mathbf{f}_{unb} \in \mathbb{C}^{K},
  \label{eq:channel}
\end{equation}
where the large-scale fading $\beta_{un}$ (path loss and shadowing) is common to the co-located antennas, and the small-scale fading is $\mathbf{f}_{unb} \sim \CN(\mathbf{0}, \mathbf{I}_K)$, independent across \glspl{oru} and \glspl{ue}. For $K = 1$, \eqref{eq:channel} is~\cite[eq.~(2)]{girycki2024twc}. We stack the channels into $\bh_{ub} = [\bh_{u1b}^T, \dots, \bh_{uNb}^T]^T \in \mathbb{C}^{M}$ and $\bH_b = [\bh_{1b}, \dots, \bh_{Ub}]^T \in \mathbb{C}^{U \times M}$; the $u$th row of $\bH_b$ is the channel of \gls{ue} $u$, and column $(n-1)K + k$ belongs to antenna $(n,k)$. The noise at \gls{ue} $u$ on \gls{prb} $b$ has power $\sigma_{ub}^2$ and comprises thermal noise and ambient external (inter-system) interference~\cite{girycki2024eucnc,girycki2024twc}.

In \gls{tdd} operation, the \gls{odu} obtains the \gls{dl} channels from \gls{ul} measurements through channel reciprocity~\cite{girycki2024twc}. While~\cite{girycki2024eucnc,girycki2024twc} assume full channel knowledge, the \gls{odu} here uses the estimate
\begin{equation}
  \bhh_{unb} = \sqrt{\beta_{un}} \Big( (1 - \epsilon)\, \mathbf{f}_{unb} + \sqrt{\epsilon (1 - \epsilon)}\, \mathbf{e}_{unb} \Big),
  \label{eq:csi}
\end{equation}
where $\mathbf{e}_{unb} \sim \CN(\mathbf{0}, \mathbf{I}_K)$ is independent of $\mathbf{f}_{unb}$ and $\epsilon \in [0, 1)$. The error $\bh_{unb} - \bhh_{unb}$ is uncorrelated with $\bhh_{unb}$ and has the energy $\epsilon\, \E\{\|\bh_{unb}\|^2\}$, so~\eqref{eq:csi} is distributed like an \gls{mmse} estimate with \gls{nmse} $\epsilon$; $\epsilon = 0$ is perfect \gls{csi}. The estimate $\bHh_b$ is stacked like $\bH_b$. All scheduling decisions below use only $\bHh_b$, whereas the \gls{se} is evaluated with the true channel $\bH_b$.

\subsection{Antenna Activation and Received Signal}\label{sec:activation}

On \gls{prb} $b$, \gls{oru} $n$ activates a subset $\mathcal{S}_{nb} \subseteq \Kset$ of its antennas, $\mathcal{S}_b = (\mathcal{S}_{1b}, \dots, \mathcal{S}_{Nb})$ collects the activations of all \glspl{oru}, and $\hat{M}_b = \sum_{n \in \Nset} |\mathcal{S}_{nb}|$ is the number of active antennas. We distinguish two granularities. At the \emph{\gls{oru} level}, an \gls{oru} transmits on \gls{prb} $b$ with all of its antennas or not at all, $\mathcal{S}_b \in \Sru = \{\mathcal{S}_b : \mathcal{S}_{nb} \in \{\emptyset, \Kset\}\ \forall n\}$; for $K = 1$ this is the per-\gls{prb} selection of serving \glspl{ap} in~\cite{girycki2024twc}. At the \emph{antenna level}, any subset is allowed, $\mathcal{S}_b \in \Sant = \{\mathcal{S}_b : \mathcal{S}_{nb} \subseteq \Kset\ \forall n\}$, and $\Sru \subseteq \Sant$. Inactive antennas do not transmit, and no \gls{fh} data are sent for them (\cref{sec:splits}). Following~\cite[eq.~(14)]{girycki2024eucnc}, the bandwidth utilization of antenna $(n,k)$ is
\begin{equation}
  \mu_{nk} = \frac{1}{B} \sum_{b \in \Bset} \ind\{k \in \mathcal{S}_{nb}\},
  \label{eq:mu}
\end{equation}
and the mean number of active antennas per \gls{prb} is $\hat{M} = \sum_{n, k} \mu_{nk} = \frac{1}{B} \sum_{b} \hat{M}_b$. When all \glspl{prb} use the same activation, $\hat{M}$ is the number of serving antennas of~\cite{girycki2024eucnc}. In NR, a \gls{ue} may assume the same precoding over a \gls{prg} of 2 or 4 \glspl{prb} or the whole bandwidth~\cite[Sec.~5.1.2.3]{ts38214}; an activation that is constant over each \gls{prg} is a special case of the per-\gls{prb} model.

\gls{oru} $n$ precodes the modulation symbols $s_{ub}$, $\E\{|s_{ub}|^2\} = 1$, of the single-layer \glspl{ue} into
\begin{equation}
  \bx_{nb} = \sum_{u \in \Uset} \bw_{nub}\, s_{ub} \in \mathbb{C}^{K}, \qquad
  [\bw_{nub}]_k = 0\ \ \forall u,\ \forall k \notin \mathcal{S}_{nb},
  \label{eq:tx}
\end{equation}
which generalizes~\cite[eq.~(2)]{girycki2024eucnc}. Let $\bxh_{nb}$ denote the signal that \gls{oru} $n$ actually radiates after the \gls{fh} transport (\cref{sec:compression}); with ideal \gls{fh}, $\bxh_{nb} = \bx_{nb}$. The signal received by \gls{ue} $u$ on \gls{prb} $b$ is
\begin{equation}
  y_{ub} = \sum_{n \in \Nset} \bh_{unb}^T \bxh_{nb} + n_{ub},
  \label{eq:rx}
\end{equation}
which reduces to~\cite[eq.~(3)]{girycki2024eucnc} for $K = 1$ and ideal \gls{fh}. With $\bw_{ub} = [\bw_{1ub}^T, \dots, \bw_{Nub}^T]^T$ and $\bW_b = [\bw_{1b}, \dots, \bw_{Ub}] \in \mathbb{C}^{M \times U}$, the effective channel $\bE_b = \bH_b \bW_b$ has the entries $[\bE_b]_{u,u'} = \sum_{n} \bh_{unb}^T \bw_{nu'b}$, the complex amplitude with which the symbol of \gls{ue} $u'$ arrives at \gls{ue} $u$; contributions of different \glspl{oru} add as complex amplitudes. As in~\cite[eq.~(5)]{girycki2024eucnc}, the sum \gls{se} is
\begin{equation}
  \SE_{\mathrm{DL}} = \frac{1}{B} \sum_{b \in \Bset} \sum_{u \in \Uset} \log_2(1 + \SINR_{ub}),
  \label{eq:se}
\end{equation}
where the \gls{sinr} with ideal \gls{fh},
\begin{equation}
  \SINR_{ub} = \frac{\big| [\bE_b]_{u,u} \big|^2}{\sum_{u' \neq u} \big| [\bE_b]_{u,u'} \big|^2 + \sigma_{ub}^2},
  \label{eq:sinr_ideal}
\end{equation}
generalizes~\cite[eq.~(4)]{girycki2024eucnc}; \cref{sec:compression} adds the compression distortion.

\subsection{Beamforming and Power}\label{sec:beamforming}

As in~\cite{girycki2024eucnc}, the beamforming type is $F \in \{\mathrm{MRT}, \mathrm{ZF}\}$, applied to all \glspl{ue}, and the precoder follows from the activation as in the two-step \gls{rrs} of~\cite{girycki2024twc}. Every active antenna is budgeted the power $\Pmax$, the maximum transmit power of an antenna per \gls{prb} in~\cite[eq.~(8)]{girycki2024eucnc}. Hence the total power on \gls{prb} $b$ and the power of each \gls{oru} are limited by
\begin{align}
  \sum_{n \in \Nset} \sum_{u \in \Uset} \|\bw_{nub}\|^2 &\le \hat{M}_b \Pmax, \label{eq:ptot}\\
  \sum_{u \in \Uset} \|\bw_{nub}\|^2 &\le K \Pmax, \quad \forall n. \label{eq:poru}
\end{align}
For $K = 1$, \eqref{eq:poru} is the per-antenna constraint of~\cite[eq.~(8)]{girycki2024eucnc} and~\cite[eq.~(10)]{girycki2024twc}, and~\eqref{eq:ptot} follows from it. For a vector $\mathbf{a} \in \mathbb{C}^K$, $\mathbf{a}_{\mathcal{S}}$ denotes $\mathbf{a}$ with its entries outside $\mathcal{S}$ set to zero.

\subsubsection{MRT}
Following Assumptions~2 and~3 of~\cite{girycki2024twc}, every active \gls{oru} serves one \gls{ue}, the one with the strongest estimated channel over its full array, $v_{nb} = \arg\max_{u} \|\bhh_{unb}\|^2$, and applies local \gls{mrt} on its active antennas:
\begin{equation}
  \bw_{n v_{nb} b} = \sqrt{|\mathcal{S}_{nb}| \Pmax}\; \frac{\bhh_{v_{nb} n b, \mathcal{S}_{nb}}^{*}}{\big\|\bhh_{v_{nb} n b, \mathcal{S}_{nb}}\big\|},
  \label{eq:mrt}
\end{equation}
and $\bw_{nub} = \mathbf{0}$ for $u \neq v_{nb}$.
Each active antenna thus contributes $\Pmax$ to the power of its \gls{oru}, and~\eqref{eq:ptot} holds with equality. For $K = 1$, \eqref{eq:mrt} is the \gls{mrt} precoder of~\cite[eqs.~(13), (14)]{girycki2024twc}. As in~\cite{girycki2024eucnc}, the \gls{mrt} weights of an \gls{oru} do not depend on the channels of other \glspl{oru} and can be computed locally at the \gls{oru}. Only activations in which every \gls{ue} is served by at least one active \gls{oru} are considered.

\subsubsection{ZF}
All active antennas serve all \glspl{ue}. With $\bHh_{b,\mathcal{S}} \in \mathbb{C}^{U \times \hat{M}_b}$ the columns of $\bHh_b$ that belong to active antennas, the precoder before normalization is
\begin{equation}
  \bWt_b = \bHh_{b,\mathcal{S}}^H \big( \bHh_{b,\mathcal{S}} \bHh_{b,\mathcal{S}}^H \big)^{-1},
  \label{eq:zf}
\end{equation}
whose rows are mapped to the active antennas; it requires $\hat{M}_b \ge U$ and $\operatorname{rank}(\bHh_{b,\mathcal{S}}) = U$. Let $q_{nb} = \sum_{k \in \mathcal{S}_{nb}} \sum_{u} |[\bWt_b]_{(n,k),u}|^2$ be the power that \gls{oru} $n$ would transmit with $\bWt_b$. Extending~\cite[eq.~(18)]{girycki2024twc}, the normalized precoder is $\bW_b = c_b \bWt_b$ with
\begin{equation}
  c_b^2 = \min\bigg\{ \frac{\hat{M}_b \Pmax}{\sum_{n} q_{nb}},\ \frac{K \Pmax}{\max_{n} q_{nb}} \bigg\}.
  \label{eq:norm}
\end{equation}
All antennas are scaled by the same factor, which preserves the nulls; the factor is set by the tighter of~\eqref{eq:ptot} and~\eqref{eq:poru}, so part of the budget generally stays unused. For $K = 1$, \eqref{eq:norm} is~\cite[eq.~(18)]{girycki2024twc}. With perfect \gls{csi}, \gls{zf} yields $\bE_b = c_b \mathbf{I}_U$. The \gls{zf} weights depend on the channels of all active antennas and are computed centrally at the \gls{odu}~\cite{girycki2024eucnc}.

\subsection{Functional Splits}\label{sec:splits}

As in~\cite{girycki2024eucnc}, two lower-layer splits are considered~\cite{tr38801,tr38816,oran_cus_v22}. They differ in where the beamforming~\eqref{eq:tx} is applied:
\begin{itemize}
  \item \emph{Split 7-1:} the \gls{odu} applies the beamforming, and the \gls{fh} carries the frequency-domain samples $[\bx_{nb}]_k$ of every active antenna $(n,k)$. In \gls{oran} terms, this is a Category-A \gls{oru} with one \gls{eaxc} per antenna.
  \item \emph{Split 7-2} (\gls{oran} 7-2x): the \gls{oru} applies the beamforming, and the \gls{fh} carries the symbols $s_{ub}$ of the \glspl{ue} it serves, $\Uset_{nb} = \{u : \bw_{nub} \neq \mathbf{0}\}$, i.e., all \glspl{ue} for \gls{zf} and $\{v_{nb}\}$ for \gls{mrt} if $\mathcal{S}_{nb} \neq \emptyset$. For \gls{zf}, the \gls{odu} also sends the weights, one block of $|\mathcal{S}_{nb}|$ weights per served \gls{ue} and \gls{prb}, in the control plane (Section Extension~11 of~\cite{oran_cus_v22}). For \gls{mrt}, the \gls{oru} computes its weights from its own \gls{ul} measurements. This is a Category-B \gls{oru}.
\end{itemize}
Following~\cite{girycki2024eucnc}, one split $s \in \{\sone, \stwo\}$ is used by all \glspl{oru} and is selected per network state, i.e., per number of \glspl{ue} and noise level. The split is a slow configuration of the \glspl{oru}, whereas the activation $\mathcal{S}_b$ changes per \gls{prb} and slot.

\subsection{Fronthaul Compression}\label{sec:compression}

The \gls{iq} payloads between \gls{odu} and \glspl{oru} are compressed with the methods of~\cite[Annex~A]{oran_cus_v22}, listed in \cref{tab:compression}. The block methods compress the 12 \glspl{resel} of one \gls{prb}, one \gls{eaxc}, and one OFDM symbol with a common parameter of $p_c$ bits and $w_c$ bits per I and Q component, which gives
\begin{equation}
  q_c = 2 w_c + p_c / 12
  \label{eq:qc}
\end{equation}
bits per \gls{resel}. Modulation compression applies only to the unprecoded constellation points of split 7-2 and represents them without loss with $Q_{m,u}$ bits per \gls{resel}, where $Q_{m,u}$ is the modulation order of the \gls{mcs} of \gls{ue} $u$~\cite{ts38214,lagen2021modcomp}. The \gls{zf} weights of split 7-2 are compressed block-wise with a weight width $w_{\mathrm{w}}$, i.e., $2 w_{\mathrm{w}} |\mathcal{S}_{nb}| + p_{\mathrm{w}}$ bits per block, where $p_{\mathrm{w}}$ counts the block parameter and the beam identifier. The uncompressed reference of~\cite{girycki2024eucnc} is $q = 32$, i.e., 16-bit fixed point. We denote by $c$ the compression configuration, i.e., the method and width of the \gls{iq} payload and, for split 7-2 with \gls{zf}, of the weights, and by $\mathcal{C}_s$ the configurations admissible for split $s$.

\begin{table}[t]
\centering
\caption{\gls{oran} compression methods~\cite[Annex~A]{oran_cus_v22} considered for the fronthaul payloads.}
\label{tab:compression}
\footnotesize
\begin{tabular}{@{}llll@{}}
\toprule
Method (udCompMeth) & Payload & Bits per \gls{resel} & Distortion \\
\midrule
Fixed point (0) & 7-1, 7-2 & $2w$ & $\rho_{\mathrm{fix}}(w)$ \\
Block floating point (1) & 7-1, 7-2, weights & $2w + 2/3$ & $\rho_{\mathrm{BFP}}(w)$ \\
Block scaling (2) & 7-1, 7-2, weights & $2w + 2/3$ & $\rho_{\mathrm{BS}}(w)$ \\
$\mu$-law (3) & 7-1, 7-2, weights & $2w + 2/3$ & $\rho_{\mu}(w)$ \\
Modulation (4) & 7-2 symbols & $Q_{m,u}$ & $0$ \\
\bottomrule
\end{tabular}
\end{table}

\emph{Distortion model.} For a lossy method, the decompressed sample is modeled as
\begin{equation}
  \mathcal{Q}_c(x) = x + e, \qquad \E\{e x^*\} = 0, \qquad \E\{|e|^2\} = \rho_c\, \E\{|x|^2\},
  \label{eq:distortion}
\end{equation}
where the relative distortion $\rho_c$ of each method and width is measured with the \gls{oran} codecs; for \gls{bfp}, $\rho_{\mathrm{BFP}}(w) \approx 3.5 \cdot 4^{-w}$, i.e., $-18.6$~dB for $w = 4$ and $6$~dB less per additional bit. \todo{Measured $\rho_c$ of all methods in \cref{tab:compression}.} The radiated signal of \gls{oru} $n$ is
\begin{equation}
  \bxh_{nb} =
  \begin{cases}
    \mathcal{Q}_c(\bx_{nb}), & s = \sone,\\[2pt]
    \sum_{u \in \Uset} \bwh_{nub}\, \mathcal{Q}_c(s_{ub}), & s = \stwo,
  \end{cases}
  \label{eq:radiated}
\end{equation}
where $\mathcal{Q}_c$ acts on each antenna stream for split 7-1 and on each \gls{ue} stream for split 7-2, and $\bwh_{nub}$ are the decompressed \gls{zf} weights or the locally computed \gls{mrt} weights. Let $\bWh_b$ collect the weights that are actually applied and $\bEh_b = \bH_b \bWh_b$; $\bEh_b = \bE_b$ except for split 7-2 with \gls{zf}. Under~\eqref{eq:distortion}, the \gls{sinr} becomes
\begin{equation}
  \SINR_{ub} = \frac{\big| [\bEh_b]_{u,u} \big|^2}{\sum_{u' \neq u} \big| [\bEh_b]_{u,u'} \big|^2 + D_{ub} + \sigma_{ub}^2},
  \label{eq:sinr}
\end{equation}
with the compression distortion
\begin{equation}
  D_{ub} =
  \begin{cases}
    \rho_c \sum_{n} \sum_{k \in \mathcal{S}_{nb}} \big| [\bh_{unb}]_k \big|^2 P_{nkb}, & s = \sone,\\[2pt]
    \rho_c \sum_{u' \in \Uset} \big| [\bEh_b]_{u,u'} \big|^2, & s = \stwo,
  \end{cases}
  \label{eq:dist}
\end{equation}
where $P_{nkb} = \sum_{u} |[\bw_{nub}]_k|^2$ is the power of antenna $(n,k)$. For split 7-1, the compression errors of the antenna streams are independent and add up non-coherently at the \gls{ue}. For split 7-2, every serving \gls{oru} receives the same compressed symbol, so its error is beamformed like the symbol itself; with modulation compression, $\rho_c = 0$. Errors of the \gls{zf} weights change $\bEh_b$: the deviation of the desired amplitude $[\bEh_b]_{u,u}$ is absorbed by the precoded \gls{dmrs}, while the leakage to other \glspl{ue} adds to the interference. With ideal \gls{fh}, $D_{ub} = 0$, $\bEh_b = \bE_b$, and~\eqref{eq:sinr} reduces to~\eqref{eq:sinr_ideal}. The \gls{se}~\eqref{eq:se} with~\eqref{eq:sinr} is denoted by $\SE_{\mathrm{DL}}(F, \{\mathcal{S}_b\}; s, c)$.

\subsection{Fronthaul Load}\label{sec:load}

Following~\cite[eq.~(9)]{girycki2024eucnc}, the basic \gls{fh} load of one stream that occupies the whole carrier with $q$ bits per \gls{resel} is
\begin{equation}
  \Gamma(q) = q\, \frac{N_{\mathrm{symb}} N_{\mathrm{sc}}}{T_{\mathrm{slot}}},
  \label{eq:gamma}
\end{equation}
where $N_{\mathrm{symb}} = 14$ is the number of OFDM symbols per slot, $N_{\mathrm{sc}} = 12B$ the number of subcarriers, and $T_{\mathrm{slot}}$ the slot duration; for 50 \glspl{prb} at 30~kHz and $q = 32$, $\Gamma = 537.6$~Mbit/s, as in~\cite{girycki2024eucnc}. Extending~\cite[eq.~(13)]{girycki2024eucnc}, the load of split 7-1 is
\begin{equation}
  R^{\sone}(c) = \sum_{n \in \Nset} \Big( \Gamma(q_c) \sum_{k \in \Kset} \mu_{nk} + \gamma_{\sone} \Big),
  \label{eq:r71}
\end{equation}
i.e., every active antenna adds a stream on the \glspl{prb} where it is active. Extending~\cite[eqs.~(10), (11)]{girycki2024eucnc}, the load of split 7-2 is
\begin{equation}
  R^{\stwo}(c) = \sum_{n \in \Nset} \Big( \sum_{u \in \Uset} \Gamma(q_{c,u})\, \nu_{nu} + R^{\mathrm{w}}_n + \gamma_{\stwo} \Big),
  \label{eq:r72}
\end{equation}
where $q_{c,u}$ is~\eqref{eq:qc} or $Q_{m,u}$ for modulation compression, and
\begin{equation}
  \nu_{nu} = \frac{1}{B} \sum_{b \in \Bset} \ind\{u \in \Uset_{nb}\}
  \label{eq:nu}
\end{equation}
is the bandwidth utilization of \gls{ue} $u$ at \gls{oru} $n$. The weights of \gls{zf}, sent once per slot and \gls{prb}, add
\begin{equation}
  R^{\mathrm{w}}_n = \frac{\ind\{F = \mathrm{ZF}\}}{T_{\mathrm{slot}}} \sum_{b \in \Bset} \sum_{u \in \Uset_{nb}} \big( 2 w_{\mathrm{w}} |\mathcal{S}_{nb}| + p_{\mathrm{w}} \big).
  \label{eq:rw}
\end{equation}
The constants $\gamma_{\stwo} = 713.9$~Mbit/s and $\gamma_{\sone} = 121$~Mbit/s are the signalling overheads of~\cite{tr38816} used in~\cite{girycki2024eucnc}. That work accounts for the beamforming weights only through $\gamma_{\stwo}$; we keep $\gamma_{\stwo}$ for comparability and make the weight load~\eqref{eq:rw} explicit because it depends on the activation and on the weight compression. For $q = 32$ and without~\eqref{eq:rw}, \eqref{eq:r71} and~\eqref{eq:r72} reduce to~\cite[eqs.~(13), (10)]{girycki2024eucnc}. As in~\cite{girycki2024eucnc}, only the \gls{dl} \gls{fh} load is considered.

\subsection{Problem Formulation}\label{sec:problem}

Let $\SE_{\max}$ be the maximum sum \gls{se} over both beamforming types and all \gls{oru}-level activations with ideal \gls{fh} and the same \gls{csi}; it is the reference of~\cite[eq.~(7)]{girycki2024eucnc} and is used for both granularities so that they are compared at the same \gls{se} target. Generalizing~\cite[eqs.~(6)--(8)]{girycki2024eucnc}, the \gls{odu} solves, for each network state and granularity $g \in \{\mathrm{RU}, \mathrm{ant}\}$,
\begin{subequations}\label{eq:problem}
\begin{align}
  \min_{s \in \{\sone, \stwo\},\ c \in \mathcal{C}_s} \quad & R^{s}\big(c;\, \hat{F}, \{\hat{\mathcal{S}}_b\}\big) \label{eq:p_obj}\\
  \text{s.t.} \quad & \big(\hat{F}, \{\hat{\mathcal{S}}_b\}\big) \in \arg\min_{F,\ \{\mathcal{S}_b \in \mathcal{S}^{g}\}} \hat{M} \label{eq:p_inner}\\
  & \text{s.t.}\ \SE_{\mathrm{DL}}(F, \{\mathcal{S}_b\}; s, c) \ge \kappa\, \SE_{\max}, \label{eq:p_se}\\
  & \phantom{\text{s.t.}}\ \eqref{eq:ptot},\ \eqref{eq:poru}, \label{eq:p_power}
\end{align}
\end{subequations}
where the precoder follows from $F$ and $\mathcal{S}_b$ through~\eqref{eq:mrt} or~\eqref{eq:zf}--\eqref{eq:norm}. The inner problem~\eqref{eq:p_inner}--\eqref{eq:p_power} selects the beamforming type and the activation that meet the \gls{se} target with the fewest active antennas, which is the \gls{rrs} problem of~\cite{girycki2024twc} with the target of~\cite{girycki2024eucnc}. The outer problem selects the split and the compression with the lowest \gls{fh} load. For $K = 1$, $g = \mathrm{RU}$, $q = 32$ without compression distortion, and $\epsilon = 0$, \eqref{eq:problem} is the problem of~\cite{girycki2024eucnc}, where $\kappa = 0.98$. The nesting follows~\cite{girycki2024eucnc}; a joint minimization of the load over all variables can only lower the optimum. Since the \gls{odu} knows only $\bHh_b$, it evaluates~\eqref{eq:p_se} with the estimated channels.

\begin{remark}[Antenna-level activation and fronthaul load]\label{rem:granularity}
For the same split, compression, and beamforming type and with perfect \gls{csi}, every \gls{oru}-level activation is also an antenna-level activation, so the optimum of~\eqref{eq:p_inner} at the antenna level uses at most as many active antennas as at the \gls{oru} level. By~\eqref{eq:r71}, the 7-1 payload decreases accordingly. The 7-2 payload~\eqref{eq:r72} depends on the \glspl{ue} served by each \gls{oru} and decreases only through the weight load~\eqref{eq:rw} of \gls{zf} and through \glspl{oru} that are switched off completely. Antenna-level activation therefore also moves the boundary between the two splits.
\end{remark}

\begin{remark}[Time scales]\label{rem:timescales}
The split $s$ and the compression configuration $c$ are selected per network state, as the split in~\cite{girycki2024eucnc}; the split is configured in the \glspl{oru}, whereas the compression parameters can be signalled per section in the user plane~\cite{oran_cus_v22}. The activation $\mathcal{S}_b$ and the precoder change per \gls{prb} (or \gls{prg}) and slot.
\end{remark}

\printbibliography

\end{document}
